\section*{Заключение}
\addcontentsline{toc}{section}{Заключение}

В результате выполнения курсовой работы было разработано веб-приложение для управления личным гардеробом с возможностью формирования пользовательских образов. Проведённый анализ предметной области показал, что цифровая организация гардероба является актуальной задачей для пользователей, которые стремятся систематизировать информацию о предметах одежды, сократить время на подбор образов и рациональнее использовать уже имеющиеся вещи.

В ходе работы была подтверждена целесообразность разработки специализированного веб-приложения, так как хранение данных о гардеробе в структурированном виде позволяет пользователю быстрее находить нужные предметы одежды, распределять их по категориям, учитывать сезонность и цвет, а также формировать готовые комплекты для повседневного использования.

На основе полученных при разработке результатов можно сформулировать следующие \textbf{выводы}:

\begin{enumerate}
\item Использование веб-технологий и фреймворка Flask позволило реализовать доступное приложение, которое не требует установки дополнительного программного обеспечения на устройство пользователя и может запускаться через обычный веб-браузер.

\item Проектирование реляционной модели данных обеспечило логичную структуру хранения информации о пользователях, категориях одежды, предметах гардероба и сохранённых образах. Реализованная структура базы данных соответствует функциональным требованиям приложения и позволяет связывать элементы одежды с конкретными пользователями и сформированными комплектами.

\item Реализация системы аутентификации, авторизации и разграничения прав доступа позволила выделить несколько типов пользователей: гостя, обычного пользователя и администратора. Такой подход обеспечивает защиту пользовательских данных и предоставляет администратору возможность управления системной информацией.

\item Разработанный CRUD-интерфейс для управления элементами гардероба обеспечивает добавление, просмотр, редактирование и удаление предметов одежды. Для каждого элемента предусмотрено указание категории, цвета, сезона и загрузка изображения, что делает гардероб наглядным и удобным для использования.

\item Реализованный модуль формирования пользовательских образов позволяет создавать комплекты одежды вручную и автоматически подбирать элементы с учётом выбранного сезона и предпочтительного цвета. Дополнительно предусмотрена возможность сохранения и выгрузки образа в виде PNG-коллажа.

\end{enumerate}

\textbf{Рекомендации и предложения.} Для практического применения разработанного решения рекомендуется использовать облачное размещение веб-приложения. Это позволит пользователю получать доступ к личному гардеробу с разных устройств без необходимости локальной установки программы. Для учебного проекта может использоваться бесплатный хостинг PythonAnywhere, поддерживающий запуск Flask-приложений и хранение данных в SQLite.

\textbf{Перспективы углубления темы.} В дальнейшем функциональность приложения может быть расширена за счёт добавления календаря планирования образов, учёта погодных условий, более сложных алгоритмов подбора комплектов, рекомендаций по сочетаемости цветов и возможности совместного использования гардероба несколькими пользователями. Также перспективным направлением является интеграция с мобильным интерфейсом, что сделает работу с цифровым гардеробом более удобной в повседневном использовании.

При разработке пояснительной записки к курсовому проекту соблюдались требования методических рекомендаций по выполнению курсовых проектов для студентов направлений 09.03.01 и 09.03.03, а также требования ГОСТ 7.32-2017, определяющие структуру и правила оформления научно-исследовательских работ~\cite{method_course_project,gost_report}.